\section{Case I Typed callback sees commit timing}\label{sec:caseI}

\subsection{Finite signature and exact ownership handoff}
This is a focused interpreted counterexample, not a general concurrency
or callback theorem. Fix one existing token identity $r$, with kinds
$U$ (uninitialized) and $S$ (usable state), and private carriers
$H_U=\{*\}$ and $H_S=\{0,1\}$. Write
\[
U=(1;r:U),\qquad E=(1;r:S),\qquad F=(\mathsf{Bit};r:S).
\]
The public/owned parts are separated by the semicolon. The expanded
endpoints are $\widehat U=1$, $\widehat E=\mathsf{Bit}$ and
$\widehat F=\mathsf{Bit}_{\rm reply}\times\mathsf{Bit}_{\rm store}$,
with their indicated token ownership. The complete primitive tables are
\begin{center}
\begin{tabular}{@{}lllll@{}}
\toprule Method&Type&Private/public action&Cost&Unary node\\\midrule
$\mathsf{init}$&$U\Rightarrow E$&$*\mapsto0$&0&$\mathsf{Init}$\\
$\mathsf{commit}$&$E\Rightarrow E$&$b\mapsto1$&1&$\mathsf{Commit}$\\
$\mathsf{read}$&$E\Rightarrow F$&$b\mapsto(b,b)$&1&$\mathsf{Read}$\\
\bottomrule
\end{tabular}
\end{center}
Each row has one witness of probability one and one output leaf owning
exactly $r$. The two rows $b=0,1$ are both admitted. Initialization only
retypes an existing token. The signature has no service-use counter:
O4 permits these finite fixed-token primitives independently of its
$\mathsf{Ready}_j$ example. Both compared wrappers perform exactly two
calls and have the identical input/output kind $S$. Thus no mismatched
counter is suppressed. The separately typed read-budget variant below
also supplies an explicit matched-counter presentation.

For each of the four maps $f:\mathsf{Bit}\to\mathsf{Bit}$, let the
rank-zero, empty-capture template be
\[
K_f:E\Rightarrow F!\{\mathsf{read}\},\qquad
K_f=\Bind(\mathsf{read},\ z\mapsto\Ret(f(z),r)).
\]
Here $r$ always denotes the currently returned owner, not a saved handle.
Let $M=\{\operatorname{code}(K_f):f:\mathsf{Bit}\to\mathsf{Bit}\}$.
The two rank-one wrapper templates capture only $p:M$ as ordinary data:
\begin{align*}
\mathsf{Early}[p]&=\Bind(\mathsf{commit},\Invoke(p)),\\
\mathsf{Late}[p]&=\Bind\bigl(\Invoke(p),\ z\mapsto
 \Bind(\operatorname{frame}_{z}(\mathsf{commit}),\Ret(z,r))\bigr).
\end{align*}
Both have type $E\Rightarrow F!\{\mathsf{commit},\mathsf{read}\}$.
The framed commit explicitly carries the ordinary bit $z$ while giving
the sole $r$ to commit; it is not a capability frame. No callback or
wrapper captures a token or private store. Each invocation strictly
decreases rank from one to zero. A common abstract implementation slot
$W:M\times E\Rightarrow F$ hides neither a public callback tag nor a
public reply, but does not give clients different inspectable wrapper
tags. Its plugging graph is acyclic.

At a callback invocation, the wrapper transfers its actual owner $r$ to
$K_f$, which transfers it to read and returns the read result's owner.
The wrapper resumes with that returned owner. During the call it retains
no executable second owner. There is no recursion, reentrant invocation,
fresh allocation or shared hidden global service.

\subsection{Actual retained interpretation}
Use the joint rig and terminal-grade causal trees of
Section~\ref{sec:core}; these particular bodies are unary trees.
For initial private bit $b$ and callback $f$, O1 occurrence composition
has one witness, mass one and total cost two. Its bodies are exactly
\begin{equation}\label{eq:I-bodies}
\begin{split}
\theta_{\mathsf{Early}}(b,f)
 &=\mathsf{Commit}(\mathsf{Read}(\leaf(f(1),1))),\\
\theta_{\mathsf{Late}}(b,f)
 &=\mathsf{Read}(\mathsf{Commit}(\leaf(f(b),1))).
\end{split}
\end{equation}
The leaf includes reply, private store and the sole $r:S$. The equations
follow by grafting the three complete primitive tables; pure relabelling,
invocation administration and ordinary framing add no event nodes.
They are actual O1 bodies, not an informal chronological trace. O3's
comparison and O4's interpretation identify these singleton families
with the corresponding finite runs.

The closing client initializes $r$, chooses $K_{\id}$ and calls $W$.
For $b=0$, both completed computations have cost two and final store one,
but their public observations are the distinct Dirac outcomes
\begin{align*}
&\mathsf{Init}(\mathsf{Commit}(\mathsf{Read}(\leaf(1,r:S)))),\\
&\mathsf{Init}(\mathsf{Read}(\mathsf{Commit}(\leaf(0,r:S)))).
\end{align*}
The deliberately coarse mathematical projection retaining only final
private store, resource endpoint and cost gives $(1,\{r:S\},2)$ in both.
This does not grant clients a private-store projection; it shows why
such an endpoint summary is inadequate. The final public reply alone
already distinguishes the programs. The difference survives a separately
declared permitted erasure of the private unary bookkeeping nodes
$\mathsf{Init},\mathsf{Commit}$.

\begin{proposition}[Interchange fails at the declared open interface]
The two wrappers are not interchangeable at the fixed causal/public-bit
observation, even though their coarse final-state summaries agree.
Neither direction has an O5 public-preserving coupling for the initialized
identity callback. At fixed expanded endpoints there is no native O1
cell between their singleton bodies.
\end{proposition}
\begin{proof}
Their only possible outcome coupling puts mass one on the pair in which
one public reply is one and the other is zero. It therefore violates
public equality, regardless of a proposed private-state relation; equal
cost does not repair the failure. The explicit closing client is admitted
by the types and ranks above. A native singleton witness map is forced,
and its probability and cost equations pass, but its body equation in
\eqref{eq:native-square} fails in~\eqref{eq:I-bodies}. After the stated
bookkeeping normalization it still fails on the public reply. Thus a
coarse endpoint equality is neither an O5 premise nor a native cell.
\end{proof}

\subsection{An explicitly matched read budget variant}
To make counters visible as well, choose a second declared signature with
kinds $S_1,S_0$ and private carrier $\mathsf{Bit}$ at each. Put
$E_j=(1;r:S_j)$ and $F_j=(\mathsf{Bit};r:S_j)$. Initialization has type
$U\Rightarrow E_1$. The sole read has type $E_1\Rightarrow F_0$, returns
the current bit and consumes the one read allowance. For $j=0,1$ declare
separate management methods
$\mathsf{commit}_j:E_j\Rightarrow E_j$, each setting the bit to one,
leaving its counter unchanged, with cost one and public label
$\mathsf{Commit}$. There is no read at $S_0$.

The rank-zero callbacks now have type $E_1\Rightarrow F_0$, and the
rank-one wrappers are
\[
\mathsf{Early}_1[p]=\mathsf{commit}_1;\Invoke(p),\qquad
\mathsf{Late}_1[p]=\Bind(\Invoke(p),\ z\mapsto
                                    \operatorname{frame}_z(\mathsf{commit}_0)).
\]
Both have exactly $E_1\Rightarrow F_0$, begin with counter one and end
with counter zero. All bodies in~\eqref{eq:I-bodies} remain the same
except for the final owned kind $S_0$; both costs remain two, and the
initialized identity callback still returns one versus zero. In
particular $\mathsf{commit}_0$ is an explicitly declared management
operation, not a forbidden read at zero. A callback cannot read again
on its $S_0$ output. This is a separate, completely specified finite
signature, not an unannounced substitution of counter types in the first
model. The same native-body and O5 obstruction proof applies.

\subsection{Exact repair and focused scope}
A valid ordering contract can require the supplied callback to be called
exactly once, after commit returns, with that returned owner, and require
the callback's returned owner to be returned to the caller. Early satisfies
this requirement; Late does not. Replacing Late by
$\Bind(\mathsf{commit},\Invoke(p))$ makes every body and cost literally
equal for all $p\in M$, giving identity native comparisons. This is an
explicit API/implementation change, not a proof that the original bodies
were equivalent. Merely requiring state-independent callback replies
would still leave different full causal trees.

The executable check \texttt{check\_case\_I\_callbacks.py} enumerates
both initial bits and the four callbacks; it verifies primitive tables,
rank decrease, actual ownership transfer, all sixteen compared wrapper
rows, costs, grafted bodies, coarse equality, distinguishing reply and
repaired equality. It also rejects invocation without the owner and a
same-rank invocation. A separate sixteen-row check covers the countered
wrappers and initialized contexts, rejects reads at $S_0$, rejects
wrong-kind management and stale handles, and verifies both exact
$E_1\Rightarrow F_0$ endpoints. Its dynamic epoch markers are ghost
ownership-audit indices, not newly allocated token identities. The proof is the complete deterministic calculation
above, not an inference from these tests. The case asserts no unbounded
callback, general transaction, serializability or concurrency result.

